\clearpage
\begin{landscape}
\begingroup
\fontsize{10}{13}\selectfont
\fancyhead[L]{\small NACHLASS. STÖRUNGEN DER PALLAS}
\renewcommand{\headrulewidth}{0pt}
\pdfbookmark[2]{Allgemeine Störungen, Erste Rechnung, [10]}{pallas-V-fragment10}
\begin{center}[10.]\end{center}
\par\noindent\hspace*{\parindent}[Zur Entwickelung der Störungen der halben grossen Axe berechnet man die Ausdrücke \(3ae\sin(v-\tilde\omega)\mathord{.}\,T\) und \(\frac{3aa\cos\varphi^2}{rr}\,V\) aus den Tafeln Seite 512—513 und entwickelt ihre Summe nach Vielfachen von \(M\); nach der Ausmultiplicirung erhält man den Ausdruck für \(\frac{\mathrm{d}\log\operatorname{hyp}n}{n\,\mathrm{d}t}\).
\par\noindent\SourceLeaf{529}\hspace*{\parindent}Sodann ist in Einheiten der 7. Decimale
\[
\frac{\mathrm{d}\log a}{n\,\mathrm{d}t}=-\tfrac{2}{3}\lambda\frac{10000000}{206265}\frac{\mathrm{d}\log\operatorname{hyp}n}{n\,\mathrm{d}t}.
\]
\par\noindent\hspace*{\parindent}Die Integration erfolgt, wie im Vorigen. Beispiel:
\begingroup
\newcommand{\VTenPoint}{\mathord{.}\mkern5mu}
\newcommand{\VTenFive}{\VTenPoint\VTenPoint\VTenPoint\VTenPoint\VTenPoint}
\begin{center}
\begin{tabular}{@{}l@{\quad}l@{}}
\(\log b\;\VTenFive\VTenFive\VTenFive\VTenPoint\VTenPoint0{,}95665\)&\\
\(\log b'\;\VTenFive\VTenFive\VTenFive\VTenPoint0{,}51851\)&\\
\multicolumn{1}{r}{\(Q=69^0\,57'\,58''\)}&\\
\(\log\sin Q\;\VTenFive\VTenFive\VTenPoint\VTenPoint\VTenPoint9{,}97289\)&\\
\(\log q\;\VTenFive\VTenFive\VTenFive\VTenPoint\VTenPoint0{,}98376\)&also das Störungsglied:\\
\(\log\left(5\frac{n'}{n}-2\right)\;\VTenFive\VTenPoint\VTenPoint\VTenPoint8{,}74966_n\)&\(-2406\cos(5M'-2M+69^0\,57'\,58'').\)\\
\cline{1-1}
\multicolumn{1}{r}{\(2{,}23410\)}&\\[3pt]
\(\log\tfrac{2}{3}\lambda\frac{10000000}{206265}\;\VTenFive\VTenPoint1{,}14727\)&\\
\cline{1-1}
\multicolumn{1}{r}{\(3{,}38137_n\)}&
\end{tabular}
\end{center}
\endgroup
\par\noindent\hspace*{\parindent}Die Störungen der halben grossen Axe ergeben sich demnach, wie folgt, wobei das säculare Glied, ebenso wie alle andern vom Argument \(M\) unabhängigen Glieder, verschwindet:
\begingroup
\renewcommand{\arraystretch}{1.15}
\begin{longtable}{@{}r@{}l@{}}
\(\delta\log a={}\)&\(+161\cos(M+93^0\,28'\,27'')\)\\
&\(+191\cos(2M+244\;6\;47)\)\\
&\(+45\cos(3M+286\;9\;45)\)\\[3pt]
&\(-40\cos(M'-4M+155^0\,24'\,31'')\)\\
&\(-99\cos(M'-3M+190\;5\;33)\)\\
&\(-146\cos(M'-2M+0\;16\;14)\)\\
&\(-64\cos(M'-M+127\;15\;14)\)\\
&\(+102\cos(M'+M+131\;53\;19)\)\\
&\(+37\cos(M'+2M+248\;38\;11)\)\\
&\(+18\cos(M'+3M+3\;0\;46)\)\\[3pt]
&\(-30\cos(2M'-5M+29^0\,29'\,10'')\)\\
&\(-38\cos(2M'-4M+16\;17\;10)\)\\
&\(-268\cos(2M'-3M+301\;0\;25)\)\\
&\(-1034\cos(2M'-2M+325\;15\;27)\)\\
&\(-2186\cos(2M'-M+135\;51\;55)\)\\
&\(+42\cos(2M'+M+180\;26\;27)\)\\
&\(+24\cos(2M'+2M+273\;44\;35)\)\\[3pt]
&\(-17\cos(3M'-6M+94^0\,34'\,43'')\)\\
&\(-20\cos(3M'-5M+230\;41\;7)\)\\
&\(-141\cos(3M'-4M+194\;14\;15)\)\\
&\(-314\cos(3M'-3M+231\;12\;2)\)\\
&\(-787\cos(3M'-2M+13\;53\;20)\)\\
&\(+1326\cos(3M'-M+145^0\,6'\,27'')\)\\
&\(+17\cos(3M'+M+190\;43\;15)\)\\[3pt]
&\(-65\cos(4M'-5M+88^0\,4'\,41'')\)\\
&\(-97\cos(4M'-4M+145\;41\;2)\)\\
&\(-325\cos(4M'-3M+272\;42\;10)\)\\
&\(-667\cos(4M'-2M+42\;6\;35)\)\\
&\(+165\cos(4M'-M+168\;53\;16)\)\\[3pt]
&\(-27\cos(5M'-6M+344^0\,7'\,4'')\)\\
&\(-41\cos(5M'-5M+76\;38\;4)\)\\
&\(-135\cos(5M'-4M+175\;2\;21)\)\\
&\(-242\cos(5M'-3M+303\;38\;18)\)\\
&\(-2406\cos(5M'-2M+69\;57\;58)\)\\
&\(+42\cos(5M'-M+192\;27\;56)\)\\[3pt]
&\(-22\cos(6M'-6M+357^0\,25'\,14'')\)\\
&\(-57\cos(6M'-5M+78\;28\;31)\)\\
&\(-116\cos(6M'-4M+204\;27\;13)\)\\
&\(-199\cos(6M'-3M+332\;22\;56)\)\\
&\(+183\cos(6M'-2M+86\;31\;33)\)\\
&\(+19\cos(6M'-M+203\;22\;3)\)\\[3pt]
&\(-24\cos(7M'-6M+355^0\,56'\,19'')\)\\
&\(-52\cos(7M'-5M+110\;12\;38)\)\\
\SourceLeaf{530}&\(-85\cos(7M'-4M+233^0\,40'\,55'')\)\\
&\(-231\cos(7M'-3M+358\;44\;52)\)\\
&\(+36\cos(7M'-2M+126\;53\;40)\)\\[3pt]
&\(-24\cos(8M'-6M+16^0\,54'\,19'')\)\\
&\(-41\cos(8M'-5M+138\;29\;50)\)\\
&\(-67\cos(8M'-4M+263\;25\;0)\)\\
&\(+305\cos(8M'-3M+27\;39\;28)\)\\[3pt]
&\(-20\cos(9M'-6M+41^0\,12'\,9'')\)\\
&\(-31\cos(9M'-5M+165\;36\;3)\)\\
&\(-62\cos(9M'-4M+291\;12\;39)\)\\
&\(+31\cos(9M'-3M+52\;59\;0)\)\\[3pt]
&\(-15\cos(10M'-6M+68^0\,51'\,43'')\)\\
&\(-23\cos(10M'-5M+191\;55\;30)\)\\
&\(-139\cos(10M'-4M+313\;17\;2)\)\\
&\(-19\cos(11M'-5M+216\;56\;27)\)\\
&\(+30\cos(11M'-4M+340\;2\;2).\)]
\end{longtable}
\endgroup

\clearpage
\endgroup
\end{landscape}
